% =============================================================================
% Capítulo 8 — Apéndices
% Conforme a APA 7 UNAD (Instructivo marzo 2023)
% =============================================================================

\section{Apéndices}

Este capítulo lista los artefactos anexos activos que respaldan el informe, la app publicada, la validación técnica y la evaluación con usuarios.

\newpage
\phantomsection
\addcontentsline{toc}{subsection}{Apéndice A: Repositorio de Código Fuente}
\label{apendice:A}
{\centering\textbf{Apéndice A}\par
\textit{Repositorio de Código Fuente}\par}

Repositorio principal del proyecto:
\begin{center}
\url{https://github.com/delarge95/WebGL-Thesis-Proposal}
\end{center}

Este repositorio, junto con el \textit{snapshot} local del proyecto y los anexos versionados, funciona como referencia técnica principal para escena, scripts, herramientas de editor y documentación de desarrollo. La copia pública de cierre se consolida en la rama \texttt{master}; por tanto, la trazabilidad pública debe verificarse contra el \texttt{HEAD} vigente de esa rama, la URL de demo publicada y los artefactos activos listados en este capítulo.

Si el repositorio público no expone todos los artefactos internos por razones de limpieza, peso o información sensible, la evaluación académica debe apoyarse en la URL de demo publicada, los anexos activos del informe, la documentación técnica interna versionada y el \textit{snapshot} local disponible para revisión.

\newpage
\phantomsection
\addcontentsline{toc}{subsection}{Apéndice B: Manual de Usuario}
\label{apendice:B}
{\centering\textbf{Apéndice B}\par
\textit{Manual de Usuario}\par}

Ruta del manual activo:
\begin{center}
\texttt{Informe\_final/Manual\_de\_usuario/manual\_usuario.tex}
\end{center}

\newpage
\phantomsection
\addcontentsline{toc}{subsection}{Apéndice C: Manual Técnico}
\label{apendice:C}
{\centering\textbf{Apéndice C}\par
\textit{Manual Técnico}\par}

Ruta del manual activo:
\begin{center}
\texttt{Informe\_final/Manual\_tecnico/manual\_tecnico.tex}
\end{center}

\newpage
\phantomsection
\addcontentsline{toc}{subsection}{Apéndice D: Documentación Técnica Activa}
\label{apendice:D}
{\centering\textbf{Apéndice D}\par
\textit{Documentación Técnica Activa}\par}

Documentos activos:

\path{Informe_final/Desarrollo_App/Documentacion_Tecnica/01_Arquitectura_del_Sistema.md}

\path{Informe_final/Desarrollo_App/Documentacion_Tecnica/02_Referencia_Tecnica_Modulos.md}

\path{Informe_final/Desarrollo_App/Documentacion_Tecnica/03_Pipeline_Renderizado_Shaders.md}

\path{Informe_final/Desarrollo_App/Documentacion_Tecnica/04_Guia_Despliegue.md}

\path{Informe_final/Desarrollo_App/Documentacion_Tecnica/05_Configuracion_WebGL.md}

\path{Informe_final/Desarrollo_App/Documentacion_Tecnica/08_Sistema_Termico_Hibrido.md}

La versión publicada para consulta desde la landing se conserva en \path{docs/manuals/}. Esta copia se sincroniza con los documentos internos cuando se actualiza el cierre técnico.

\newpage
\phantomsection
\addcontentsline{toc}{subsection}{Apéndice E: Matriz de Desconexiones}
\label{apendice:E}
{\centering\textbf{Apéndice E}\par
\textit{Matriz de Desconexiones}\par}

La matriz maestra de desconexiones entre app y documentación se encuentra en:
\begin{center}
\path{Informe_final/Desarrollo_App/MATRIZ_DESCONEXIONES_APP_DOCUMENTACION.md}
\end{center}

\newpage
\phantomsection
\addcontentsline{toc}{subsection}{Apéndice F: Catálogo Canónico del Holybro X500 V2}
\label{apendice:F}
{\centering\textbf{Apéndice F}\par
\textit{Catálogo Canónico del Holybro X500 V2}\par}

La taxonomía semántica del proyecto se fija a partir de:
\begin{center}
\path{desarrollo/unity_project/Assets/Resources/x500v2_info_panel_bilingual_catalog.json}
\end{center}

Esta fuente pública integra los identificadores canónicos, las fichas bilingües de información y las claves de trazabilidad utilizadas por el panel de la app. La taxonomía académica conserva 28 piezas canónicas; la escena final opera adicionalmente con 30 \textit{anchors} y 257 \textit{renderers}/\textit{colliders} auditados.

\newpage
\phantomsection
\addcontentsline{toc}{subsection}{Apéndice G: Validación Funcional de la Escena Final}
\label{apendice:G}
{\centering\textbf{Apéndice G}\par
\textit{Validación Funcional de la Escena Final}\par}

El documento de referencia para la build funcional es:
\begin{center}
\path{Informe_final/Desarrollo_App/VALIDACION_FUNCIONAL_FINA_2026-04-09.md}
\end{center}

\newpage
\phantomsection
\addcontentsline{toc}{subsection}{Apéndice H: Instrumentos y Registros de Evaluación}
\label{apendice:H}
{\centering\textbf{Apéndice H}\par
\textit{Instrumentos y Registros de Evaluación}\par}

\setcounter{table}{0}
\renewcommand{\thetable}{H\arabic{table}}
\setcounter{figure}{0}
\renewcommand{\thefigure}{H\arabic{figure}}

Los instrumentos metodológicos activos del cierre son el cuestionario SUS, el NASA-TLX Raw, el protocolo \textit{Think-Aloud}, la guía de mediciones técnicas WebGL, la guía de tareas comparativas 3D vs.\ 2D y las tablas de KPIs técnicos asociadas a sesiones fechadas del profiler interno.

Las tareas de prueba se alinean con cuatro acciones equivalentes: ubicar un componente, interpretar su función, explicar una relación espacial y resolver una acción de análisis visual guiado. En la condición 3D estas tareas se apoyan en navegación, selección, \textit{bottom sheet}, \textit{Explode}, \texttt{Cut} y modos visuales; en la condición 2D se resuelven mediante el paquete documental estático de referencia.

Para mantener coherencia con la metodología vigente, SUS se reporta como lectura global del prototipo 3D y NASA-TLX Raw se reporta como comparación de carga de trabajo percibida entre la condición 3D y el soporte 2D de referencia.

\begin{table}[htbp]
    \caption{Matriz de Equivalencia Informativa entre Visor 3D y Soporte 2D}
    \label{tab:equivalencia_soporte_2d_visor_3d}
    \begin{tabular}{p{3.2cm}p{4.4cm}p{4.4cm}}
        \hline
        \textbf{Tarea} & \textbf{Condición 3D} & \textbf{Condición 2D} \\ \hline
        Ubicación & Navegación orbit/pan/zoom, selección y foco. & Vista general, lámina o página del soporte. \\
        Interpretación & Ficha inferior de pieza y categoría. & Ficha textual o fragmento documental 2D. \\
        Relación espacial & Inspección 3D, aislamiento, explode o corte. & Despiece o diagrama estático. \\
        Análisis guiado & Uso de explode, cut, X-Ray, thermal o studio. & Interpretación de vista estática equivalente. \\ \hline
    \end{tabular}
    \apanote{La equivalencia es funcional, no formal: el visor 3D ofrece interacción espacial, mientras que el soporte 2D conserva información suficiente para resolver la misma intención de tarea. Elaboración propia.}
\end{table}

Las versiones activas y versionadas de instrumentos y registros se conservan en:

\path{Informe_final/validacion/CUESTIONARIO_SUS.md}

\path{Informe_final/validacion/CUESTIONARIO_NASA_TLX.md}

\path{Informe_final/validacion/GUIA_TAREAS_VALIDACION.md}

\path{Informe_final/validacion/PROTOCOLO_THINK_ALOUD.md}

\path{Informe_final/validacion/06_GUIA_MEDICIONES_TECNICAS_WEBGL.md}

\path{Informe_final/validacion/07_TABLAS_RENDIMIENTO_WEBGL_MEDICIONES.tex}

Los registros completos por participante se conservan como evidencia interna de evaluación y no forman parte de la superficie pública del repositorio. En el capítulo 5 se reportan únicamente datos agregados y participantes anonimizados; si un jurado o asesor requiere auditar los registros individuales, estos deben entregarse por un canal académico controlado y no mediante publicación pública.

Como respaldo visual del proceso de medición técnica, la Figura \ref{fig:evidencia_capturas_profiler_interno} muestra una captura puntual del profiler interno conservada junto con las capturas de cierre de la aplicación.

\begin{figure}[htbp]
    \caption{Muestra de una Captura del Profiler Interno}
    \label{fig:evidencia_capturas_profiler_interno}
    \includegraphics[width=0.92\textwidth,height=0.45\textheight,keepaspectratio]{figures/screenshots_contextual/fig_profiler_internal_evidence.png}
    \apanote{La figura documenta una captura puntual del profiler interno integrada como evidencia visual en \path{Informe_final/figures/screenshots_contextual/fig_profiler_internal_evidence.png}. Su función es respaldar visualmente el proceso de recolección técnica; los valores oficiales de FPS, memoria, renderers, mallas y triángulos reportados en las tablas técnicas provienen de las exportaciones JSON/CSV del profiler interno, no de lectura manual de esta imagen. Elaboración propia.}
\end{figure}

\newpage
\phantomsection
\addcontentsline{toc}{subsection}{Apéndice I: Estado del Pipeline de Activos}
\label{apendice:I}
{\centering\textbf{Apéndice I}\par
\textit{Estado del Pipeline de Activos}\par}

\setcounter{table}{0}
\renewcommand{\thetable}{I\arabic{table}}

Para efectos de cierre documental:

Blender y Unity conforman la ruta defendible de la versión final activa.

El FBX runtime final se documenta como unión de masters e instancias; su trazabilidad exige reporte de importación Unity, revisión de grupos, fasteners, hotspots, filtros y hélices.

Cualquier ruta nativa de LOD se documenta como decisión evaluada; si existen herramientas o scripts posteriores de LOD, no se presentan como mejora validada del estudio mientras no exista profiling y QA comparativo sobre la build publicada.

Houdini se mantiene como herramienta evaluada o experimental, no como pipeline oficial de la build cerrada, salvo evidencia nueva derivada del reimport final.

Las texturas/materiales no publicados deben mantenerse como mejora futura o documentarse expresamente si se integran a una build posterior.

El estado actual del pipeline puede verse consolidado en la Tabla \ref{tab:pipeline_3d_cierre}.

\begin{table}[htbp]
    \caption{Estado del Pipeline 3D al Cierre Documental}
    \label{tab:pipeline_3d_cierre}
    \begin{tabular}{p{2.8cm}p{2.8cm}p{2.2cm}p{4.2cm}}
        \hline
        \textbf{Etapa} & \textbf{Herramienta} & \textbf{Estado} & \textbf{Lectura editorial} \\ \hline
        Limpieza & Blender & Activa & Ruta defendible del cierre para ajuste de malla. \\
        Conversión CAD & STEPper / MoI3D & Activa / evaluada & Combinación de apertura rápida y teselación. \\
        Materiales & Unity URP & Activa & Flujo material visible en la build final. \\
        Integración & Unity & Activa & Núcleo del producto final y punto de verdad. \\
        Auditoría & Binder runtime & Activa & Coherencia entre escena y documentación. \\
        Nativa LOD & Investigación & Evaluada & Evolución de planeación técnica. \\
        Houdini & Evaluación & No canónico & Citado como exploración o trabajo futuro. \\ \hline
    \end{tabular}
    \apanote{Esta tabla evita presentar como pipeline oficial herramientas que no quedaron integradas de forma cerrada en la build final. Elaboración propia.}
\end{table}

\newpage
\phantomsection
\addcontentsline{toc}{subsection}{Apéndice J: Demo Público}
\label{apendice:J}
{\centering\textbf{Apéndice J}\par
\textit{Demo Público}\par}

La URL pública del demo publicado para la entrega es:
\begin{center}
\url{https://delarge95.github.io/WebGL-Thesis-Proposal/}
\end{center}

Esta URL corresponde al punto de acceso a la build WebGL publicada. Las métricas del capítulo 5 registran la fecha y el contexto de medición utilizado para evitar inconsistencias entre documento y producto.

\newpage
\phantomsection
\addcontentsline{toc}{subsection}{Apéndice K: Audio y Alcance No Integrado}
\label{apendice:K}
{\centering\textbf{Apéndice K}\par
\textit{Audio y Alcance No Integrado}\par}

Al momento de este cierre documental:

No existen clips de audio integrados en \texttt{Assets}.

Por defecto, el audio se clasifica como trabajo futuro si no entra en la build publicada.

El capítulo 5 reporta la evaluación técnica y de usuarios con los datos disponibles para esta versión del informe.

\newpage
\phantomsection
\addcontentsline{toc}{subsection}{Apéndice L: Auditoría Matemática y Trazabilidad con Wolfram}
\label{apendice:L}
{\centering\textbf{Apéndice L}\par
\textit{Auditoría Matemática y Trazabilidad con Wolfram}\par}

La auditoría activa de coherencia entre fórmulas documentadas, lógica de código y arquitectura real se consolida en:
\begin{center}
\path{Informe_final/Desarrollo_App/AUDITORIA_MATEMATICA_Y_ARQUITECTURA_2026-04-12.md}
\end{center}

Como soporte técnico adicional, el subsistema térmico se documenta en el anexo técnico público:
\begin{center}
\path{Informe_final/Desarrollo_App/Documentacion_Tecnica/08_Sistema_Termico_Hibrido.md}
\end{center}

Ambos documentos respaldan la interpretación correcta del modo \textit{Thermal} como simulación híbrida heurística y no como FEA en tiempo real. Los registros históricos usados durante el desarrollo permanecen como evidencia interna cuando no forman parte de la superficie pública del repositorio.

\newpage
\phantomsection
\addcontentsline{toc}{subsection}{Apéndice M: Correlación entre la Propuesta y el Informe Final}
\label{apendice:M}
{\centering\textbf{Apéndice M}\par
\textit{Correlación entre la Propuesta y el Informe Final}\par}

\setcounter{table}{0}
\renewcommand{\thetable}{M\arabic{table}}

La siguiente tabla resume la correlación de las tablas presentadas en la propuesta (véase la Tabla \ref{tab:correlacion_tablas_propuesta_informe}).

\begin{table}[htbp]
    \caption{Correlación entre Tablas de la Propuesta y el Informe Final}
    \label{tab:correlacion_tablas_propuesta_informe}
    \begin{tabular}{p{3.8cm}p{3.4cm}p{4.8cm}}
        \hline
        \textbf{Tabla en la propuesta} & \textbf{Estado en informe} & \textbf{Justificación} \\ \hline
        Comparativa Web 3D & Capítulo 2 & Sustento del estado del arte y selección de Unity. \\
        Cronograma general & Solo en propuesta & Artefacto de planeación; pierde valor en informe final. \\
        Presupuesto estimado & Solo en propuesta & Viabilidad ex ante; no es un artefacto técnico. \\
        Resultados esperados & Tabla entregados vs. esperados & Requiere contrastar expectativa inicial con entrega real. \\ \hline
    \end{tabular}
    \apanote{La ausencia de una tabla del documento final no implica incoherencia si su función era exclusivamente prospectiva en la propuesta. Elaboración propia.}
\end{table}

Asimismo, la Tabla \ref{tab:productos_entregados_esperados} presenta un contraste entre los productos entregados y los originalmente esperados.

\begin{table}[htbp]
    \caption{Productos Entregados vs.\ Esperados al Cierre Documental}
    \label{tab:productos_entregados_esperados}
    \begin{tabular}{p{3.8cm}p{2.4cm}p{5.8cm}}
        \hline
        \textbf{Producto esperado} & \textbf{Estado} & \textbf{Observación de cierre} \\ \hline
        Software web funcional & Cerrado & Flujo funcional publicado y documentable; métricas agregadas. \\
        Shaders URP optimizados & Entregado & Modos visibles/ocultos documentados de forma diferenciada. \\
        Modelos 3D optimizados & Entregado & Pipeline defendible y escena funcional para MVP. \\
        Documento de tesis & En cierre & Capítulos saneados con datos reales de rendimiento y usuarios. \\
        Informe de usabilidad & Entregado & Capítulo 5 consolida SUS, NASA-TLX Raw y Think-Aloud. \\ \hline
    \end{tabular}
    \apanote{Esta tabla actualiza el horizonte de productos de la propuesta sin borrar la diferencia entre expectativa inicial y estado real de cierre. Elaboración propia.}
\end{table}

\newpage
\phantomsection
\addcontentsline{toc}{subsection}{Apéndice N: Resumen de Artefactos Activos del Cierre}
\label{apendice:N}
{\centering\textbf{Apéndice N}\par
\textit{Resumen de Artefactos Activos del Cierre}\par}

\setcounter{table}{0}
\renewcommand{\thetable}{N\arabic{table}}

A continuación se presentan los artefactos activos resultantes del cierre (véase la Tabla \ref{tab:resumen_artefactos_cierre}).

\begin{table}[htbp]
    \caption{Resumen de Artefactos Activos del Cierre Documental}
    \label{tab:resumen_artefactos_cierre}
    \begin{tabular}{p{3.5cm}p{3.5cm}p{1.8cm}p{3.2cm}}
        \hline
        \textbf{Artefacto} & \textbf{Ubicación principal} & \textbf{Estado} & \textbf{Uso en defensa} \\ \hline
        Informe final actual & \path{Informe_final/informe_final.pdf} & Activo & Evaluación académica principal. \\
        Manual de usuario & \path{Informe_final/Manual_de_usuario/manual_usuario.pdf} & Activo & Soporte al flujo visible. \\
        Manual técnico & \path{Informe_final/Manual_tecnico/manual_tecnico.pdf} & Activo & Soporte de arquitectura. \\
        Landing y docs & \path{docs/index.html}, \path{docs/manuals/} & Activo & Punto de acceso a build y docs. \\
        Validación e inst. & \path{Informe_final/validacion/} & Activo & Evidencia pública anonimizada. \\
        Telemetría agregada & \path{Telemetria/Mediciones_WebGL/} & Activo & Exportaciones profiler. \\
        Diagramas Mermaid & \path{Informe_final/figures/chapter4/} & Activo & Explican arquitectura y flujo. \\
        Anexos técnicos & \path{Informe_final/Desarrollo_App/Documentacion_Tecnica/} & Activo & Documentan WebGL y térmica. \\ \hline
    \end{tabular}
    \apanote{Esta tabla resume únicamente artefactos activos para consulta pública del cierre; auditorías históricas, bitácoras de trabajo, borradores y registros individuales se conservan como evidencia interna o por canal académico controlado, pero no como fuente primaria pública. Elaboración propia.}
\end{table}
